Pemetaan tanaman pertanian berperan penting dalam mendukung ketahanan pangan, perencanaan pola tanam, serta pengelolaan sumber daya agrikultur. 
Citra satelit Sentinel-2 dengan resolusi tinggi spasial dan temporal telah banyak dimanfaatkan dalam klasifikasi tutupan lahan, termasuk untuk identifikasi jenis tanaman. 
Penelitian ini mengadopsi pendekatan \textit{semantic segmentation} berbasis \textit{deep learning}, yaitu DeepLabV3+ dan SegFormer, untuk mengklasifikasikan tanaman pertanian berdasarkan citra multi-temporal Sentinel-2 dan label dari \textit{Cropland Data Layer} (CDL) milik USDA. 

Studi literatur menunjukkan bahwa pendekatan multi-temporal dan kombinasi band spektral dapat meningkatkan akurasi klasifikasi karena mampu menangkap dinamika spektral dan fenologi tanaman. 
Area studi terletak di Sacramento Valley, California, yang merupakan salah satu kawasan agrikultur paling produktif di Amerika Serikat. 
Data Sentinel-2 dikumpulkan setiap 15 hari selama periode 2022--2024 dan diproses melalui tahapan \textit{masking} awan, \textit{histogram matching}, \textit{resampling} spasial, normalisasi, dan \textit{patching}. 

Evaluasi dilakukan menggunakan metrik \textit{Overall Accuracy}, \textit{Mean Intersection over Union} (mIoU), dan F1-score, serta dilengkapi analisis kesalahan dan waktu inferensi. 
Penelitian ini juga mengimplementasikan \textit{attention mechanism} pada masing-masing model untuk meningkatkan performa segmentasi. 
Eksperimen dilakukan dalam beberapa skenario yang menguji pengaruh data multi-temporal dan kombinasi spektral terhadap hasil klasifikasi tanaman. 
Hasil dari studi ini diharapkan dapat memberikan kontribusi terhadap pengembangan model segmentasi citra multispektral yang lebih akurat dan adaptif terhadap variasi spasial dan temporal tutupan lahan pertanian.

\vspace{0.5cm}
\textbf{Kata Kunci:} Pemetaan Pertanian, Multi-temporal, Band Spektral, Semantic Segmentation
